\subsection{常系数线性方程组}

沿用 \(n^{\circ}\) 8 的记号，我们更一般地考虑形如下的由 m 个微分方程组成的方程组

\[
\sum_ {k = 1} ^ {n} p _ {j k} (\mathrm{D}) x _ {k} = b _ {j} (t) \quad (1 \leqslant j \leqslant m)\tag{47}
\]

的方程组，其中未知函数 \(x_{k}\) ( \(1 \leqslant k \leqslant n\) ) 与右端 \(b_{j}\) ( \(1 \leqslant j \leqslant m\) ) 都是实变量 t 的复值函数，而 \(p_{jk}(D)\) 是微分算子 D 的常（复）系数多项式（任意次数）（ \(1 \leqslant j \leqslant m, 1 \leqslant k \leqslant n\) ）。

这样的方程组与 IV, p.~1164（公式 (5)）中所考虑的那些不同型，如下例所示：

\[
\left\{ \begin{array}{c} \mathrm{D} x _ {1} = a (t) \\ \mathrm{D} ^ {2} x _ {1} + \mathrm{D} x _ {2} + x _ {2} = b (t). \end{array} \right.\tag{48}
\]

我们将把自己限于函数 \(b_{j}(t)\) 在区间 J 上无穷次可导的情形，并且只寻求在 J 上无穷次可导的解 \((x_{k})_{1\leqslant k\leqslant n}\) 。令 \(\mathbf{b}(t)=\left(b_{1}(t),\ldots,b_{m}(t)\right)\) （这是由 J 到 \(C^{m}\) 中的一个映射）以及 \(\mathbf{x}=\left(x_{1},x_{2},\ldots,x_{n}\right)\) ，方程组 (47) 便可写成

\[
P (\mathrm{D}) \mathbf {x} = \mathbf {b} (t)\tag{49}
\]

其中 \(P(\mathrm{D})\) 是矩阵 \((p_{jk}(\mathrm{D}))\) ，有 m 行 n 列，其系数属于 D 的多项式、以 C 为系数的环 C{[}D{]}。设 \(f_{j}(\mathrm{D})\) ( \(1 \leqslant j \leqslant r \leqslant \operatorname{Min}(m, n)\) ) 是矩阵 \(P(\mathrm{D})\) 的非零相似不变量；我们知道 (Alg., VII, p.~32)，它们是一些完全确定的首一多项式，使得 \(f_{j}\) 整除 \(f_{j+1}\) （对 \(1 \leqslant j \leqslant r - 1\) ；r 是 \(P(\mathrm{D})\) 的秩）；此外，存在两个方阵 \(U(\mathrm{D})\) 与 \(V(\mathrm{D})\) ，阶数分别为 m 与 n，它们是可逆的（在阶数分别为 m 与 n、以 D 的复系数多项式环 C{[}D{]} 的元素为系数的方阵环中），并且使得矩阵 \(Q(\mathrm{D}) = (q_{jk}(\mathrm{D})) = U(\mathrm{D}) P(\mathrm{D}) V(\mathrm{D})\) 的所有元素除对角元 \(q_{jj}(\mathrm{D}) = f_{j}(\mathrm{D})\) （对 \(1 \leqslant j \leqslant r\) ）外都为零。现在令 \(\mathbf{y} = V^{-1}(\mathbf{D})\mathbf{x}\) ；方程 (49) 等价于方程

\[
U (\mathrm{D}) \big (P (\mathrm{D}) \big (V (\mathrm{D}) \mathbf {y} \big) \big) = U (\mathrm{D}) \mathbf {b}
\]

也就是

\[
Q (\mathrm{D}) \mathbf {y} = U (\mathrm{D}) \mathbf {b}\tag{50}
\]

因为 \(U(\mathbf{D})\) 可逆。现在，若 \(\mathbf{y} = (y_1, y_2, \ldots, y_n)\) ，且若

\[
U (\mathrm{D}) \mathbf {b} (t) = \left(c _ {1} (t), \dots , c _ {m} (t)\right),
\]

则方程 (50) 可以写成

\[
f _ {j} (\mathrm{D}) y _ {j} = c _ {j} (t) \quad \text {   for   } 1 \leqslant j \leqslant r\tag{51}
\]

\[
0 = c _ {j} (t) \quad \text {   for   } r + 1 \leqslant j \leqslant m.\tag{52}
\]

因此，除非条件 (52) 满足，该方程组没有任何无穷次可导的解；这时 \(y_{j}\) （下标 \(j \leqslant r\) ）的确定就化为对 r 个常系数线性微分方程 (51) 的积分；而 \(y_{j}\) （下标 \textgreater{} r 的那些）是任意的无穷次可导函数。这样确定了 (50) 的解 y 之后，就由公式 \(\mathbf{x} = V(\mathrm{D})\mathbf{y}\) 得出 (47) 的解。

注. 1) 某些多项式 \(f_{j}(\mathrm{D})\) 可以化为非零常数；这时相应的 \(y_{j}\) 就被完全确定。

\begin{enumerate}
\def\labelenumi{\arabic{enumi})}
\setcounter{enumi}{1}
\item
  当 \(b_{j}\) 全为零时，即当方程组 (47) 是齐次的时，条件 (52) 总是满足的；进一步，若 \(r = n\) ，则可看出 (47) 的解所成之集是 \(\mathbf{C}\) 上的一个向量空间，其维数等于诸 \(f_{j}(\mathrm{D})\) 的次数之和，也就是 \(\det(P(\mathrm{D}))\) 的次数。
\item
  给定多项式 \(p_{jk}(\mathbf{D})\) ，一个当右端无穷次可导（或某个阶数可导）时容许解的方程组 (47)，在右端是任意规则函数时可能不容许解：例 (48) 表明了这一点，当 \(a(t)\) 没有原函数时它没有解。我们在这里不去寻求当右端是任意规则函数时出现的补充可解性条件。
\end{enumerate}

\setcounter{exercisegroup}{0}
\refstepcounter{exercisegroup}
